% ── §4 Information-Theoretic Measures ──
\section{Information-Theoretic Measures}
\label{sec:theory-info}

\subsection{Shannon Entropy}
\label{sec:shannon}

For an alignment column $i$ containing $N$ sequences, the Shannon entropy
\cite{shannon1948} quantifies the evolutionary variability at that position:
\begin{equation}
	H(i) = -\sum_{a=1}^{20} P_i(a) \log_2 P_i(a),
	\label{eq:shannon}
\end{equation}
where $P_i(a)$ denotes the observed frequency of amino acid $a$ at column $i$,
computed after excluding gap characters (state~20).  The entropy ranges from
$H = 0$ (fully conserved column) to $H = \log_2(20) \approx 4.32$ bits
(maximally variable).  We classify a column as \emph{variable} if $H(i)
	> 0.3$, a threshold chosen empirically to capture positions where at least two
or three different amino acids appear with appreciable frequency.

\subsection{Mutual Information}
\label{sec:mutual-info}

For a pair of alignment columns $(i, j)$, the mutual information measures the
statistical dependence between the amino-acid distributions at those columns:
\begin{equation}
	\MI(i,j) = \sum_{a=1}^{20}\sum_{b=1}^{20} P_{ij}(a,b)
	\log_2 \frac{P_{ij}(a,b)}{P_i(a)\,P_j(b)},
	\label{eq:mi}
\end{equation}
where $P_{ij}(a,b)$ is the joint frequency of observing residue $a$ at column
$i$ and residue $b$ at column $j$.  MI is measured in bits and equals zero
when the two positions are statistically independent.  We compute two variants:
the \emph{full MI}, which counts all observations including the reference
(majority-residue) pair, and the \emph{mutation-only MI}, which excludes the
reference pair to isolate mutation-driven covariation from conservation-driven
signal.

Raw MI suffers from two well-known biases: phylogenetic inertia inflates
correlations between sequences sharing recent common ancestry, and random
sampling noise produces nonzero baseline values even for independent
positions.  The Average Product Correction (APC) of Dunn et
al.\cite{dunn2008} removes much of this bias by subtracting the product of row
and column means from the MI matrix:
\begin{equation}
	\mathrm{MI}_{\mathrm{APC}}(i,j) =
	\MI(i,j) - \frac{\bar{\MI}(i,\cdot)\,\bar{\MI}(\cdot,j)}{\bar{\MI}},
	\label{eq:apc}
\end{equation}
where $\bar{\MI}(i,\cdot)$ is the mean MI over all partners of position $i$
and $\bar{\MI}$ is the grand mean.

\subsection{Perplexity Ratio}
\label{sec:perplexity}

While MI measures total statistical dependence, it does not distinguish between
many weakly contributing states and a few strongly determined ones.  To capture
this distinction, we define the \emph{perplexity ratio} for a pair $(i,j)$ as
\begin{equation}
	r(i,j) = \frac{\mathrm{PP}(j)}
	{\frac{1}{|\mathcal{A}|}\sum_{a \in \mathcal{A}}
		\mathrm{PP}(j \mid i = a)},
	\label{eq:pp-ratio}
\end{equation}
where $\mathrm{PP}(j) = 2^{H(j)}$ is the marginal perplexity of column $j$,
$\mathrm{PP}(j \mid i = a) = 2^{H(j|i=a)}$ is the conditional perplexity given
that column $i$ holds residue $a$, and $\mathcal{A}$ is the set of residues at
column $i$ with at least five observations.  A ratio exceeding unity indicates
that knowing the residue at position $i$ reduces the effective number of
choices at position $j$.  In our PSICOV150 evaluation, the perplexity ratio
achieved the highest Spearman correlation with native-contact fraction among
all sequence-derived metrics tested ($\rho = 0.670$, $p = 0.006$).

\subsection{Coupling Constants}
\label{sec:coupling}

The coupling constant for a residue pair $(a, b)$ at positions $(i, j)$ is
defined as the log-odds ratio
\begin{equation}
	J_{ij}(a,b) = \ln \frac{P_{ij}(a,b)}{P_i(a)\,P_j(b)},
	\label{eq:coupling}
\end{equation}
which equals zero when the pair occurs at its independent (marginal-product)
frequency.  Positive values indicate coevolutionary enrichment; negative
values indicate anti-correlation or negative selection.  The sigmoid transform
$\sigma(J) = (1 + e^{-J})^{-1}$ maps couplings to probabilities in $[0, 1]$.
